% ---------------------------------------------------------------------------
% Solar Array Rooftop Survey — group activity presentation (advanced version)
%
% Build:  latexmk -pdf main.tex      (or: pdflatex main.tex, twice)
% Notes:  uncomment the \setbeameroption line below for a speaker-notes build.
% ---------------------------------------------------------------------------
\documentclass[11pt]{beamer}

% ---------------------------------------------------------------------------
% preamble.tex — theme, colours, and reusable TikZ styles for the deck.
% Kept separate from main.tex so the slides themselves stay readable.
% ---------------------------------------------------------------------------

\usepackage[T1]{fontenc}
\usepackage[utf8]{inputenc}
\usepackage{lmodern}
\usepackage{amsmath}
\usepackage{booktabs}
\usepackage{tikz}
\usetikzlibrary{angles,quotes,calc,arrows.meta,positioning,patterns,decorations.pathreplacing}

% --- Palette --------------------------------------------------------------
\definecolor{solarnavy}{RGB}{26,42,66}     % headings, structure
\definecolor{solarslate}{RGB}{74,90,112}   % secondary text
\definecolor{solaramber}{RGB}{233,157,42}  % the "solar" accent
\definecolor{solarteal}{RGB}{31,133,136}   % second accent (bearings)
\definecolor{solarmist}{RGB}{240,243,247}  % panel backgrounds
\definecolor{solarrule}{RGB}{206,214,224}  % hairlines

% --- Beamer theme ---------------------------------------------------------
\usetheme{default}
\usecolortheme{default}
\setbeamertemplate{navigation symbols}{}
\setbeamercolor{normal text}{fg=solarnavy}
\setbeamercolor{structure}{fg=solarteal}
\setbeamercolor{frametitle}{fg=solarnavy}
\setbeamercolor{framesubtitle}{fg=solarslate}
\setbeamercolor{alerted text}{fg=solaramber}
\setbeamerfont{frametitle}{size=\Large,series=\bfseries}
\setbeamerfont{framesubtitle}{size=\small,series=\mdseries}
\setbeamerfont{title}{size=\huge,series=\bfseries}

\setbeamertemplate{itemize item}{\color{solaramber}\rule[0.35ex]{0.5em}{0.5em}}
\setbeamertemplate{itemize subitem}{\color{solarteal}\rule[0.35ex]{0.35em}{0.35em}}
\setbeamertemplate{enumerate item}{\color{solarteal}\bfseries\insertenumlabel.}
\setbeamercolor{itemize/enumerate body}{fg=solarnavy}

\makeatletter
\setbeamertemplate{frametitle}{%
  \nointerlineskip
  \begin{beamercolorbox}[wd=\paperwidth,sep=0pt,leftskip=0.9cm,rightskip=0.9cm]{frametitle}
    \vskip0.35cm
    {\usebeamerfont{frametitle}\insertframetitle\par}%
    \ifx\insertframesubtitle\@empty\else
      {\vskip1pt\usebeamerfont{framesubtitle}\usebeamercolor[fg]{framesubtitle}%
       \insertframesubtitle\par}%
    \fi
    \vskip4pt{\color{solaramber}\rule{2.4em}{2pt}}\par\vskip0.05cm
  \end{beamercolorbox}%
}
\makeatother

\setbeamersize{text margin left=0.9cm, text margin right=0.9cm}

\setbeamertemplate{footline}{%
  \hbox{}\hfill
  \begin{beamercolorbox}[wd=\paperwidth,ht=2.4ex,dp=1.6ex,leftskip=0.9cm,rightskip=0.9cm]{}
    {\color{solarslate}\scriptsize\insertshorttitle}\hfill
    {\color{solarslate}\scriptsize\insertframenumber}
  \end{beamercolorbox}%
}

% --- Reusable bits --------------------------------------------------------

% A soft callout panel: \panel{content}
\newcommand{\panel}[1]{%
  \begingroup\setlength{\fboxsep}{8pt}%
  \colorbox{solarmist}{\parbox{\dimexpr\textwidth-16pt\relax}{#1}}%
  \endgroup}

% Same, but flagged with an amber edge for the "read this twice" boxes
\newcommand{\keypanel}[1]{%
  \begingroup\setlength{\fboxsep}{8pt}%
  \colorbox{solaramber!18}{\parbox{\dimexpr\textwidth-16pt\relax}{#1}}%
  \endgroup}

% Inline emphasis in the accent colours. Mode-aware so the same macro can be
% used inside an align* body as well as in running text.
\newcommand{\amber}[1]{{\color{solaramber}\ifmmode\mathbf{#1}\else\bfseries #1\fi}}
\newcommand{\teal}[1]{{\color{solarteal}\ifmmode\mathbf{#1}\else\bfseries #1\fi}}

% --- Shared TikZ styles ---------------------------------------------------
\tikzset{
  beam/.style   = {solarnavy, line width=1.4pt, line cap=round},
  ray/.style    = {solarteal, line width=1.1pt, line cap=round},
  north/.style  = {solarslate, line width=0.8pt, dashed, -{Stealth[length=4pt]}},
  helper/.style = {solarrule, line width=0.7pt},
  vertex/.style = {circle, fill=solarnavy, inner sep=1.5pt},
  vlabel/.style = {solarnavy, font=\small\bfseries},
  dim/.style    = {solarslate, font=\scriptsize},
  angmark/.style= {solaramber, line width=1pt},
  panelcell/.style = {fill=solarnavy!85, draw=solarmist, line width=0.3pt},
}

% Small reusable north arrow: \northarrow{<length in cm>}
\newcommand{\northarrow}[1]{%
  \draw[north] (0,0) -- (90:#1) node[above, dim] {N};}

% ---------------------------------------------------------------------------
% diagrams.tex — every TikZ figure in the deck, one macro each.
%
% Geometry note used throughout: a bearing is measured clockwise from north,
% while TikZ measures anticlockwise from east, so
%
%     tikz angle = 90 - bearing
%
% Bearing 038 -> 52 degrees, bearing 142 -> -52 degrees. The two beams
% therefore sit symmetrically either side of the east axis, which keeps every
% figure in this deck consistent with the survey's own orientation.
% ---------------------------------------------------------------------------

% === 1. The site: a triangular corner of a warehouse rooftop ===============
\newcommand{\diagRooftop}{%
\begin{tikzpicture}[scale=1]
  % roof slab + parapet
  \fill[figpale] (0,0) rectangle (9.6,5.6);
  \draw[figrule, line width=2pt] (0,0) rectangle (9.6,5.6);
  \draw[figrule, line width=0.6pt] (0.18,0.18) rectangle (9.42,5.42);

  % rooftop plant, to sell "this is a real roof, not a textbook"
  \foreach \x/\y/\w/\h in {6.4/3.7/1.5/1.2, 8.1/1.1/1.0/1.3, 6.6/1.2/0.9/0.7}{
    \fill[figmuted!25] (\x,\y) rectangle ++(\w,\h);
    \draw[figmuted!60, line width=0.6pt] (\x,\y) rectangle ++(\w,\h);
  }
  \draw[figmuted!45, line width=2.5pt] (7.15,3.7) -- (7.15,2.6) -- (8.6,2.6) -- (8.6,2.4);
  \node[dim, anchor=west] at (6.4,5.05) {existing HVAC plant};

  % north arrow
  \begin{scope}[shift={(0.85,4.35)}]
    \draw[figmuted, line width=0.9pt, -{Stealth[length=5pt]}] (0,0) -- (0,0.85);
    \node[dim, above] at (0,0.85) {N};
  \end{scope}

  % the survey triangle
  \coordinate (O) at (2.3,2.6);
  \coordinate (A) at ($(O)+(52:3.045)$);
  \coordinate (B) at ($(O)+(-52:2.394)$);
  \fill[figaccent!25] (O) -- (A) -- (B) -- cycle;
  \draw[figaccent!55, line width=0.8pt] (A) -- (B);
  \draw[beam] (A) -- (O) -- (B);

  \foreach \p in {O,A,B}{\node[vertex] at (\p) {};}
  \node[vlabel, left=3pt]        at (O) {O};
  \node[vlabel, above right=1pt] at (A) {A};
  \node[vlabel, below right=1pt] at (B) {B};

  % side labels ride the outside of each beam, clear of the vertex at O
  \node[dim, rotate=52,  anchor=south, yshift=1pt] at ($(O)!0.62!(A)$) {beam OA, 34.6 m};
  \node[dim, rotate=-52, anchor=north, yshift=-1pt] at ($(O)!0.62!(B)$) {beam OB, 27.2 m};
  \node[figink, font=\scriptsize\bfseries, align=center]
    at ($(O)!0.72!($(A)!0.5!(B)$)$) {unused\\corner};

  \node[dim, anchor=south west] at (0.25,-0.55) {Plan view of the warehouse rooftop, bounded by two structural beams};
\end{tikzpicture}}

% === 2. What a bearing actually is ========================================
\newcommand{\diagBearingConvention}{%
\begin{tikzpicture}[scale=1]
  \foreach \cx/\mth/\lab/\name in {0/52/038/A, 4.4/-52/142/B}{
    \begin{scope}[shift={(\cx,0)}]
      \draw[figrule, line width=0.8pt] (0,0) circle (1.45);
      \foreach \a/\t in {90/N, 0/E, -90/S, 180/W}{
        \draw[figrule, line width=0.8pt] (\a:1.32) -- (\a:1.45);
        \node[dim] at (\a:1.72) {\t};
      }
      \draw[north] (0,0) -- (90:1.32);
      \draw[ray, -{Stealth[length=5pt]}] (0,0) -- (\mth:1.32);
      \node[vertex] at (0,0) {};
      \node[vlabel, below left=1pt] at (0,0) {O};
      \node[vlabel] at (\mth:1.72) {\name};
      % clockwise sweep from north
      \draw[angmark, -{Stealth[length=4pt]}] (90:0.72) arc[start angle=90, end angle=\mth, radius=0.72];
      \node[figalert, font=\small\bfseries] at ({(90+\mth)/2}:1.05) {$\lab^\circ$};
    \end{scope}}
  \node[dim, align=center, text width=12cm] at (2.2,-2.35)
    {A bearing is always measured \emph{clockwise from north}, from the same origin each time.};
\end{tikzpicture}}

% === 3. Subtracting the bearings to get the included angle ================
\newcommand{\diagBearingSubtraction}{%
\begin{tikzpicture}[scale=1]
  \coordinate (O) at (0,0);
  \coordinate (A) at (52:3.9);
  \coordinate (B) at (-52:3.07);

  \draw[north] (O) -- (90:3.3);
  \node[dim, above] at (90:3.3) {N};

  \fill[figalert!14] (O) -- (52:1.05) arc[start angle=52, end angle=-52, radius=1.05] -- cycle;

  \draw[beam] (O) -- (A);
  \draw[beam] (O) -- (B);
  \node[vertex] at (O) {}; \node[vertex] at (A) {}; \node[vertex] at (B) {};
  \node[vlabel, below left=1pt]  at (O) {O};
  \node[vlabel, above right=1pt] at (A) {A};
  \node[vlabel, below right=1pt] at (B) {B};
  \node[dim, rotate=52,  anchor=south] at ($(O)!0.85!(A)$) {34.6 m};
  \node[dim, rotate=-52, anchor=north] at ($(O)!0.85!(B)$) {27.2 m};

  % bearing arcs, nested at different radii so they never collide
  \draw[ray, -{Stealth[length=4pt]}] (90:1.35) arc[start angle=90, end angle=52, radius=1.35];
  \node[figaccent, font=\small\bfseries] at (71:1.66) {$038^\circ$};
  \draw[ray, -{Stealth[length=4pt]}] (90:2.15) arc[start angle=90, end angle=-52, radius=2.15];
  \node[figaccent, font=\small\bfseries] at (14:2.48) {$142^\circ$};

  % the included angle: the thing we are actually after
  \draw[angmark] (52:1.05) arc[start angle=52, end angle=-52, radius=1.05];
  \node[figalert, font=\bfseries] at (0:0.72) {$104^\circ$};
\end{tikzpicture}}

% === 4. The general rule, including the case that needs adjusting =========
\newcommand{\diagBearingRule}{%
\begin{tikzpicture}[scale=0.95]
  % --- case 1: difference under 180, take it as it is
  \begin{scope}
    \draw[figrule, line width=0.8pt] (0,0) circle (1.5);
    \draw[north] (0,0) -- (90:1.5); \node[dim, above] at (90:1.5) {N};
    \fill[figalert!14] (0,0) -- (52:0.85) arc[start angle=52, end angle=-52, radius=0.85] -- cycle;
    \draw[ray, -{Stealth[length=4pt]}] (0,0) -- (52:1.5);
    \draw[ray, -{Stealth[length=4pt]}] (0,0) -- (-52:1.5);
    \draw[angmark] (52:0.85) arc[start angle=52, end angle=-52, radius=0.85];
    \node[figalert, font=\small\bfseries] at (0:1.22) {$104^\circ$};
    \node[dim] at (52:1.78) {$038^\circ$};
    \node[dim] at (-52:1.78) {$142^\circ$};
    \node[figink, font=\small, align=center] at (0,-2.35)
      {$142-38 = 104 < 180$\\[2pt] take it directly};
  \end{scope}
  % --- case 2: difference over 180, take the reflex complement
  \begin{scope}[shift={(6.4,0)}]
    \draw[figrule, line width=0.8pt] (0,0) circle (1.5);
    \draw[north] (0,0) -- (90:1.5); \node[dim, above] at (90:1.5) {N};
    % bearings 040 and 300 -> tikz 50 and -210 (=150)
    \fill[figalert!14] (0,0) -- (50:0.85) arc[start angle=50, end angle=150, radius=0.85] -- cycle;
    \draw[ray, -{Stealth[length=4pt]}] (0,0) -- (50:1.5);
    \draw[ray, -{Stealth[length=4pt]}] (0,0) -- (150:1.5);
    \draw[angmark] (50:0.85) arc[start angle=50, end angle=150, radius=0.85];
    \node[figalert, font=\small\bfseries] at (100:1.15) {$100^\circ$};
    \node[dim] at (50:1.8) {$040^\circ$};
    \node[dim] at (150:1.85) {$300^\circ$};
    \node[figink, font=\small, align=center] at (0,-2.35)
      {$300-40 = 260 > 180$\\[2pt] use $360 - 260 = 100^\circ$};
  \end{scope}
\end{tikzpicture}}

% === 5. The oblique triangle, reduced to SAS ==============================
\newcommand{\diagSAS}{%
\begin{tikzpicture}[scale=1]
  \coordinate (O) at (0,0);
  \coordinate (A) at (52:3.29);
  \coordinate (B) at (-52:2.58);

  \fill[figaccent!12] (O) -- (A) -- (B) -- cycle;
  \draw[beam] (O) -- (A);
  \draw[beam] (O) -- (B);
  \draw[figmuted, line width=1pt, dashed] (A) -- (B);

  \draw[angmark] (52:0.8) arc[start angle=52, end angle=-52, radius=0.8];
  \node[figalert, font=\bfseries] at (0:1.14) {$104^\circ$};
  \node[dim, anchor=east] at (0:0.42) {$C$};

  \foreach \p in {O,A,B}{\node[vertex] at (\p) {};}
  \node[vlabel, left=3pt]        at (O) {O};
  \node[vlabel, above right=1pt] at (A) {A};
  \node[vlabel, below right=1pt] at (B) {B};

  \node[figink, font=\small\bfseries, rotate=52,  anchor=south] at ($(O)!0.58!(A)$) {$b = 34.6$ m};
  \node[figink, font=\small\bfseries, rotate=-52, anchor=north] at ($(O)!0.58!(B)$) {$a = 27.2$ m};
  \node[dim, anchor=west, align=left, text width=2.6cm] at ($(A)!0.5!(B)+(0.3,0)$)
    {side AB, never measured and never needed};
\end{tikzpicture}}

% === 6. Where sin C comes from: it is the height you never measured =======
\newcommand{\diagHeight}{%
\begin{tikzpicture}[scale=1]
  \coordinate (O) at (0,0);
  \coordinate (B) at (2.86,0);
  \coordinate (A) at (104:3.63);
  \coordinate (F) at (-0.879,0);

  \fill[figaccent!12] (O) -- (A) -- (B) -- cycle;
  \draw[helper, dashed] (F) -- (O);          % base line, extended past O
  \draw[beam] (A) -- (O) -- (B);
  \draw[figmuted, line width=1pt, dashed] (A) -- (B);

  % the perpendicular height
  \draw[figaccent, line width=1.2pt, dashed] (A) -- (F);
  \draw[figaccent, line width=0.8pt] ($(F)+(0.22,0)$) -- ($(F)+(0.22,0.22)$) -- ($(F)+(0,0.22)$);
  \node[figaccent, font=\small\bfseries, rotate=90, anchor=south] at ($(A)!0.5!(F)$)
    {$h = b\sin C$};

  % angles at O
  \draw[angmark] (0:0.8) arc[start angle=0, end angle=104, radius=0.8];
  \node[figalert, font=\small\bfseries] at (52:1.12) {$104^\circ$};
  \draw[figaccent, line width=1pt] (104:0.42) arc[start angle=104, end angle=180, radius=0.42];
  \node[figaccent, font=\scriptsize\bfseries] at (142:0.68) {$76^\circ$};

  \foreach \p in {O,A,B}{\node[vertex] at (\p) {};}
  \node[vlabel, below right=1pt] at (O) {O};
  \node[vlabel, above right=1pt] at (A) {A};
  \node[vlabel, below right=1pt] at (B) {B};
  \node[dim, below=3pt] at ($(O)!0.55!(B)$) {$a = 27.2$ m};
  % right-aligned so the label runs away from O rather than over it
  \node[dim, anchor=north east] at ($(F)+(0.18,-0.14)$) {foot of the height};

  % Two separate nodes rather than one: TikZ's align=left redefines \\, and a
  % \\[dim] nested inside a {...} group within the node text breaks it.
  \node[figink, anchor=north west, font=\small, align=left] at (3.5,3.35)
    {Area $=\tfrac12 \times$ base $\times$ height\\[3pt]
     $\phantom{\text{Area}}=\tfrac12\, a\,(b\sin C)$};
  \node[figmuted, anchor=north west, font=\scriptsize, align=left] at (3.5,2.3)
    {$C$ is obtuse, so the foot lands\\[1pt]
     outside the triangle, and\\[1pt]
     $\sin 104^\circ = \sin 76^\circ$ still holds.};
\end{tikzpicture}}

% === 7. Illustrative panel layout (schematic, not an engineering layout) ==
\newcommand{\diagPanels}{%
\begin{tikzpicture}[scale=1]
  \coordinate (O) at (0,0);
  \coordinate (A) at (52:3.29);
  \coordinate (B) at (-52:2.58);

  \fill[figpale] (O) -- (A) -- (B) -- cycle;
  \begin{scope}
    \clip (O) -- (A) -- (B) -- cycle;
    \foreach \i in {0,...,8}{
      \foreach \j in {0,...,21}{
        \fill[panelcell] ({0.12+\i*0.26},{-2.3+\j*0.25}) rectangle ++(0.21,0.19);
      }}
  \end{scope}
  \draw[beam] (A) -- (O) -- (B);
  \draw[figmuted, line width=1pt] (A) -- (B);
  \foreach \p in {O,A,B}{\node[vertex] at (\p) {};}
  \node[vlabel, left=2pt]        at (O) {O};
  \node[vlabel, above right=1pt] at (A) {A};
  \node[vlabel, below right=1pt] at (B) {B};
  \node[dim, align=center, anchor=north] at (1.6,-3.0) {schematic only, ignoring tilt, row\\spacing and shading};
\end{tikzpicture}}


% \setbeameroption{show notes}            % speaker notes interleaved
% \setbeameroption{show notes on second screen=right}

\title[Rooftop Solar Array]{Sizing a Rooftop Solar Array}
\subtitle{Finding the area of an oblique triangle from survey bearings}
\author[Group MON]{Group Monday Righteous}
\institute[]{Trigonometry Group Activity}
\date{\today}   % TODO: set the presentation date

\begin{document}

\begin{frame}
  \titlepage
\end{frame}

\begin{frame}{Outline}
  \tableofcontents
\end{frame}

% ===========================================================================
\section{The problem}

\begin{frame}{The brief}
  An engineering firm is scoping a small solar array in an unused triangular
  corner of a warehouse rooftop, boxed in by two existing structural beams.
  They need \alert{the enclosed area}, which is what sets how many panels the
  site can hold.

  \begin{block}{The catch}
    The surveyor never recorded an angle. The field notes contain
    \structure{two bearings}, which is how rooftop and land surveys are
    actually reported.
  \end{block}

  So the job splits in two: recover the angle, then find the area.
  \note{Open on the practical question, not the formula. The audience should
  want the area before we hand them a method for finding it.}
\end{frame}

\begin{frame}{The site}
  \vfill
  \begin{center}
    \scalebox{0.85}{\diagRooftop}
  \end{center}
  \vfill
  \note{Point out that the two beams are fixed, so the building gives the
  shape and nobody chose it. That is why it is an awkward oblique triangle
  instead of a convenient right one.}
\end{frame}

% ===========================================================================
\section{From bearings to an angle}

\begin{frame}{What the surveyor actually recorded}
  % widths deliberately sum to less than 1 so the block and the text beside it
  % are not shoulder to shoulder
  \begin{columns}[T,onlytextwidth]
    \begin{column}{0.48\textwidth}
      \begin{block}{Field notes at station O}
        \small
        \begin{tabular}{@{}ll@{}}
          Beam OA & 34.6 m, bearing $038^\circ$\\
          Beam OB & 27.2 m, bearing $142^\circ$\\[2pt]
          \midrule
          Angle at O & \emph{not recorded}\\
        \end{tabular}
      \end{block}
    \end{column}
    \hfill
    \begin{column}{0.44\textwidth}
      \small A total station measures \structure{distance and direction}, and
      direction is logged as a bearing. Nobody stands in the corner with a
      protractor, so \alert{the included angle has to be derived.}
    \end{column}
  \end{columns}
  \vfill
  \begin{center}
    \scalebox{0.70}{\diagBearingConvention}
  \end{center}
  \note{Budget a beat here. The audience needs to accept that "no angle given"
  is normal survey practice, not a trick in the question.}
\end{frame}

\begin{frame}{From two bearings to one angle}
  \begin{columns}[T]
    \begin{column}{0.46\textwidth}
      \small
      Two conditions make this a plain subtraction:
      \begin{itemize}
        \item both bearings share the \structure{same reference point}, O;
        \item both are measured \structure{clockwise from north}.
      \end{itemize}
      \medskip

      \begin{alertblock}{Included angle}
        \centering $\angle AOB = 142^\circ - 38^\circ = 104^\circ$
      \end{alertblock}

      No trigonometry yet. This step is bookkeeping about how the directions
      were recorded.
    \end{column}
    \begin{column}{0.54\textwidth}
      \centering
      \scalebox{0.80}{\diagBearingSubtraction}
    \end{column}
  \end{columns}
  \note{This is the slide that makes the scenario real rather than
  plug-and-chug. Do not rush it.}
\end{frame}

\begin{frame}{The rule in general}
  \vfill
  \begin{center}
    \scalebox{0.75}{\diagBearingRule}
  \end{center}
  \vfill
  {\small Subtract the bearings; if the result exceeds $180^\circ$ you have
  measured the \emph{reflex} side, so take $360^\circ$ minus it instead. Here
  $104^\circ < 180^\circ$, so no adjustment is needed.}
  \note{Include this even though our case does not need it. It shows the
  method was applied knowingly rather than by luck.}
\end{frame}

% ===========================================================================
\section{Finding the area}

\begin{frame}{Now it is a standard SAS triangle}
  \begin{columns}[T]
    \begin{column}{0.44\textwidth}
      \begin{block}{Known}
        \begin{tabular}{@{}ll@{}}
          $a$ & $=$ OB $= 27.2$ m\\[2pt]
          $b$ & $=$ OA $= 34.6$ m\\[2pt]
          $C$ & $= 104^\circ$ (derived)\\
        \end{tabular}
      \end{block}
      \small
      The angle sits \alert{between} the two known sides, which is exactly the
      SAS pattern.
      \medskip

      $C$ is obtuse, so the triangle is genuinely oblique: no right angle to
      fall back on, and no perpendicular height in the field notes.
    \end{column}
    \begin{column}{0.56\textwidth}
      \centering
      \scalebox{0.80}{\diagSAS}
    \end{column}
  \end{columns}
\end{frame}

\begin{frame}{Why $\tfrac12ab\sin C$ is the right tool}
  \vfill
  \begin{center}
    \scalebox{0.80}{\diagHeight}
  \end{center}
  \vfill
  {\small The familiar $\tfrac12\times\text{base}\times\text{height}$ never
  went away. \alert{$\sin C$ simply supplies the height} from the side and
  angle we do have.}
  \note{The one line worth saying out loud: sin C substitutes for a
  perpendicular height that was never directly measured.}
\end{frame}

\begin{frame}{The calculation}
  \begin{columns}[T]
    \begin{column}{0.55\textwidth}
      \begin{align*}
        \text{Area} &= \tfrac12\,ab\sin C\\[3pt]
                    &= \tfrac12 \times 34.6 \times 27.2 \times \sin 104^\circ\\[3pt]
                    &= 470.56 \times \sin 104^\circ\\[3pt]
                    &= 470.56 \times 0.9703\\[3pt]
                    &\approx 456.58\ \text{m}^2
      \end{align*}
    \end{column}
    \begin{column}{0.45\textwidth}
      \begin{block}{Watch for}
        \small
        Calculator in \structure{degrees}, not radians.
        \medskip

        $\sin 104^\circ = \sin 76^\circ = 0.9703$. An obtuse angle gives a
        perfectly ordinary positive sine.
        \medskip

        Round only at the \structure{final step}.
      \end{block}
    \end{column}
  \end{columns}
\end{frame}

\begin{frame}{Complete solution}
  \begin{center}
    \small
    \begin{tabular}{@{}ll@{}}
      \toprule
      \textbf{Quantity} & \textbf{Value}\\
      \midrule
      Side OA ($b$)               & 34.6 m\\
      Side OB ($a$)               & 27.2 m\\
      Bearing to A                & $038^\circ$\\
      Bearing to B                & $142^\circ$\\
      Included angle $C$          & $142^\circ - 38^\circ = 104^\circ$\\
      $\sin 104^\circ$            & $\approx 0.9703$\\
      $\tfrac12 \times$ OA $\times$ OB & 470.56\\
      \midrule
      \textbf{Area}               & $\mathbf{\approx 456.58\ \text{m}^2}$\\
      \bottomrule
    \end{tabular}
  \end{center}
\end{frame}

\begin{frame}{The answer}
  \begin{alertblock}{Area of the triangular rooftop section}
    \begin{equation*}
      \text{Area} = \tfrac12 \times 34.6 \times 27.2 \times \sin 104^\circ
                  \approx 456.58\ \text{m}^2
    \end{equation*}
  \end{alertblock}
  \vfill
  \begin{center}
    From two beam lengths and two bearings, with no angle ever measured.
  \end{center}
  \vfill
\end{frame}

% ===========================================================================
\section{Discussion}

\begin{frame}{A note on rounding}
  \begin{columns}[T]
    \begin{column}{0.5\textwidth}
      \small
      $\sin 104^\circ = 0.970296\ldots$, so
      \[470.56 \times 0.970296 = 456.582\ldots\]

      A rounded $0.9709$ would give $456.87\ \text{m}^2$, about
      $0.3\ \text{m}^2$ adrift, roughly \alert{one panel's worth} of area.
      \medskip

      Small, but the kind of error that compounds: carry the full value
      through and round once, at the end.
    \end{column}
    \begin{column}{0.5\textwidth}
      \begin{block}{Sensitivity of the answer}
        \small
        \begin{tabular}{@{}ll@{}}
          \midrule
          $\sin C = 0.9703$ & $456.58\ \text{m}^2$\\
          $\sin C = 0.9709$ & $456.87\ \text{m}^2$\\
          \midrule
          \multicolumn{2}{@{}l}{difference: $0.29\ \text{m}^2$}\\
        \end{tabular}
      \end{block}
    \end{column}
  \end{columns}
  \note{Optional slide. Drop it if time is short, but it is cheap credibility
  if a marker checks the arithmetic.}
\end{frame}

\begin{frame}{Optional: what could that area produce?}
  \begin{columns}[T]
    \begin{column}{0.56\textwidth}
      \small
      Utility-scale arrays are often quoted at roughly $170\ \text{W}$ of
      installed capacity per square metre:
      \begin{equation*}
        456.58 \times 170 \approx 77\,619\ \text{W} \approx 77.6\ \text{kW}
      \end{equation*}

      \begin{alertblock}{Say this out loud}
        \small A rough, order-of-magnitude estimate, not an engineering
        calculation. It ignores tilt, row spacing, shading and inverter
        losses, all of which reduce the real figure.
      \end{alertblock}
    \end{column}
    \begin{column}{0.44\textwidth}
      \centering
      \scalebox{0.72}{\diagPanels}
    \end{column}
  \end{columns}
  \note{Not required by the rubric. If you use it, label it as an estimate
  before you say the number, not after.}
\end{frame}

\begin{frame}{What to take away}
  \begin{block}{1. The angle was earned, not given}
    Bearings share a reference direction and a rotational sense, so their
    difference \emph{is} the angle between the lines. Real survey data arrives
    this way; converting it is part of the problem.
  \end{block}

  \begin{block}{2. It is still just base $\times$ height}
    $\tfrac12ab\sin C$ is not a new formula to memorise. $\sin C$ stands in
    for a perpendicular height that was never measured, which is why SAS is
    enough.
  \end{block}

  \vfill
  \begin{center}
    From \structure{two distances and two directions} to
    \alert{$456.58\ \text{m}^2$}, and a panel count.
  \end{center}
\end{frame}

\begin{frame}
  \vfill
  \begin{center}
    \Large That's all!
  \end{center}
  \vfill
\end{frame}

\end{document}
